%%%
%%% Radical "⿃"
%%%
\section*{Radical 196: ``⿃'' (鸟)}\addcontentsline{toc}{section}{Radical 196: ⿃,鸟}\addcontentsline{loh}{figure}{\#\#\#\# 196: ⿃}

%%%%%%%%%% 鸟 %%%%%%%%%%
\subsection*{鸟}\addcontentsline{loh}{figure}{鸟}

\begin{Entry}{鸟}{5}{⿃}[Kangxi 196]
  \begin{Phonetics}{鸟}{diao3}
    \definition{s.}{(em romances tradicionais, como termo pejorativo) maldito; condenado; fudido; o mesmo que 屌}
  \seealsoref{屌}{diao3}
  \end{Phonetics}
  \begin{Phonetics}{鸟}{niao3}[][HSK 2]
    \definition*{s.}{Sobrenome: Niao}
    \definition[只,群]{s.}{pássaro; ave}
  \end{Phonetics}
\end{Entry}

\begin{Entry}{鸟儿}{5,2}{⿃,⼉}
  \begin{Phonetics}{鸟儿}{niao3r5}
    \definition[只]{s.}{pássaro; ave}
  \end{Phonetics}
\end{Entry}

\begin{Entry}{鸟巢}{5,11}{⿃,⼮}
  \begin{Phonetics}{鸟巢}{niao3chao2}[][HSK 7-9]
    \definition[个,些]{s.}{ninho de pássaro; ninhos usados por pássaros para pôr ovos e incubar seus filhotes}
  \end{Phonetics}
\end{Entry}

%%%%%%%%%% 鸡 %%%%%%%%%%
\subsection*{鸡}\addcontentsline{loh}{figure}{鸡}

\begin{Entry}{鸡}{7}{⿃}
  \begin{Phonetics}{鸡}{ji1}[][HSK 2]
    \definition*{s.}{Sobrenome: Ji}
    \definition[只]{s.}{galo, galinha, frango | palavra ofensiva para uma mulher que ganha dinheiro fazendo sexo com homens}
  \end{Phonetics}
\end{Entry}

\begin{Entry}{鸡西}{7,6}{⿃,⾑}
  \begin{Phonetics}{鸡西}{ji1xi1}
    \definition*{s.}{Cidade no nível da prefeitura de Jixi, na província de Heilongjiang (黑龙江), no nordeste da China}
  \seealsoref{黑龙江}{hei1long2jiang1}
  \end{Phonetics}
\end{Entry}

\begin{Entry}{鸡蛋}{7,11}{⿃,⾍}
  \begin{Phonetics}{鸡蛋}{ji1dan4}[][HSK 1]
    \definition[个,枚,筐,箱,打]{s.}{ovo de galinha}
  \end{Phonetics}
\end{Entry}

%%%%%%%%%% 鸣 %%%%%%%%%%
\subsection*{鸣}\addcontentsline{loh}{figure}{鸣}

\begin{Entry}{鸣}{8}{⿃}
  \begin{Phonetics}{鸣}{ming2}
    \definition*{s.}{Sobrenome: Ming}
    \definition{v.}{(pássaros, animais ou insetos) gritar; chamar | produzir um som; tocar; fazer com que se produza um som | expressar; verbalizar; transmitir pelo ar}
  \end{Phonetics}
\end{Entry}

%%%%%%%%%% 鸦 %%%%%%%%%%
\subsection*{鸦}\addcontentsline{loh}{figure}{鸦}

\begin{Entry}{鸦}{9}{⿃}
  \begin{Phonetics}{鸦}{ya1}
    \definition[只]{s.}{corvo}
  \end{Phonetics}
\end{Entry}

\begin{Entry}{鸦雀无声}{9,11,4,7}{⿃,⾫,⽆,⼠}
  \begin{Phonetics}{鸦雀无声}{ya1que4-wu2sheng1}[][HSK 7-9]
    \definition{expr.}{``Nem mesmo um corvo ou pardal pode ser ouvido.''; silêncio absoluto; nenhum som podia ser ouvido; descreve um estado de extremo silêncio; também pode ser uma metáfora para uma situação em que ninguém se atreve a falar ou todos permanecem em silêncio}
  \synonymref{无话可说}{wu2hua4-ke3shuo1}
  \antonymref{沸沸扬扬}{fei4fei4yang2yang2}
  \antonymref{七嘴八舌}{qi1zui3-ba1she2}
  \end{Phonetics}
\end{Entry}

%%%%%%%%%% 鸭 %%%%%%%%%%
\subsection*{鸭}\addcontentsline{loh}{figure}{鸭}

\begin{Entry}{鸭}{10}{⿃}
  \begin{Phonetics}{鸭}{ya1}
    \definition[只]{s.}{pato | Gíria: prostituto}
  \end{Phonetics}
\end{Entry}

\begin{Entry}{鸭子}{10,3}{⿃,⼦}
  \begin{Phonetics}{鸭子}{ya1zi5}[][HSK 5]
    \definition[只,群]{s.}{pato | Gíria: prostituto}
  \end{Phonetics}
\end{Entry}

%%%%%%%%%% 鸵 %%%%%%%%%%
\subsection*{鸵}\addcontentsline{loh}{figure}{鸵}

\begin{Entry}{鸵}{10}{⿃}
  \begin{Phonetics}{鸵}{tuo2}
    \definition[只]{s.}{avestruz}
  \end{Phonetics}
\end{Entry}

\begin{Entry}{鸵鸟}{10,5}{⿃,⿃}
  \begin{Phonetics}{鸵鸟}{tuo2niao3}
    \definition{s.}{avestruz}
  \end{Phonetics}
\end{Entry}

%%%%%%%%%% 鸽 %%%%%%%%%%
\subsection*{鸽}\addcontentsline{loh}{figure}{鸽}

\begin{Entry}{鸽}{11}{⿃}
  \begin{Phonetics}{鸽}{ge1}
    \definition[只]{s.}{pombo}[和平鸽。===Pomba da Paz.]
  \end{Phonetics}
\end{Entry}

\begin{Entry}{鸽子}{11,3}{⿃,⼦}
  \begin{Phonetics}{鸽子}{ge1zi5}[][HSK 7-9]
    \definition[只,对,群]{s.}{pombo}
  \end{Phonetics}
\end{Entry}

%%%%%%%%%% 鹅 %%%%%%%%%%
\subsection*{鹅}\addcontentsline{loh}{figure}{鹅}

\begin{Entry}{鹅}{12}{⿃}
  \begin{Phonetics}{鹅}{e2}[][HSK 7-9]
    \definition[只,群]{s.}{ganso}
  \end{Phonetics}
\end{Entry}

%%%%%%%%%% 鹏 %%%%%%%%%%
\subsection*{鹏}\addcontentsline{loh}{figure}{鹏}

\begin{Entry}{鹏}{13}{⿃}
  \begin{Phonetics}{鹏}{peng2}
    \definition*{s.}{Peng, um pássaro gigante lendário da mitologia chinesa}
  \end{Phonetics}
\end{Entry}

\begin{Entry}{鹏程万里}{13,12,3,7}{⿃,⽲,⼀,⾥}
  \begin{Phonetics}{鹏程万里}{peng2cheng2-wan4li3}[][HSK 7-9]
    \definition{suf.}{``Um futuro brilhante o aguarda!'';  ``O lendário Roca voa dez mil milhas.''; tenha um futuro brilhante; as perspectivas futuras de alguém são brilhantes}
  \end{Phonetics}
\end{Entry}

%%%%%%%%%% 鹤 %%%%%%%%%%
\subsection*{鹤}\addcontentsline{loh}{figure}{鹤}

\begin{Entry}{鹤}{15}{⿃}
  \begin{Phonetics}{鹤}{he4}
    \definition*{s.}{Sobrenome: He}
    \definition[只]{s.}{grou (ave)}
  \end{Phonetics}
\end{Entry}

\begin{Entry}{鹤立鸡群}{15,5,7,13}{⿃,⽴,⿃,⽺}
  \begin{Phonetics}{鹤立鸡群}{he4li4ji1qun2}[][HSK 7-9]
    \definition{expr.}{destaque"-se da multidão; manifestamente superior; muito acima do comum; como um guindaste em pé entre galinhas, fique de pé acima dos outros}
  \end{Phonetics}
\end{Entry}

%%%%%%%%%% 鹦 %%%%%%%%%%
\subsection*{鹦}\addcontentsline{loh}{figure}{鹦}

\begin{Entry}{鹦}{16}{⿃}
  \begin{Phonetics}{鹦}{ying1}
    \definition[只]{s.}{papagaio}
  \end{Phonetics}
\end{Entry}

\begin{Entry}{鹦鹉}{16,13}{⿃,⿃}
  \begin{Phonetics}{鹦鹉}{ying1wu3}
    \definition[只,群]{s.}{papagaio (ave)}
  \end{Phonetics}
\end{Entry}

%%%%%%%%%% 鹰 %%%%%%%%%%
\subsection*{鹰}\addcontentsline{loh}{figure}{鹰}

\begin{Entry}{鹰}{18}{⿃}
  \begin{Phonetics}{鹰}{ying1}[][HSK 7-9]
    \definition[只,群]{s.}{falcão; águia}
  \end{Phonetics}
\end{Entry}

%%%%%%%%%% 鷄 %%%%%%%%%%
\subsection*{鷄}\addcontentsline{loh}{figure}{鷄}

\begin{Entry}{鷄}{21}{⿃}
  \begin{Phonetics}{鷄}{ji1}
    \variantof{鸡}
  \end{Phonetics}
\end{Entry}

%%%%% EOF %%%%%

